\chapter*{Introduction to Part~I}
\addcontentsline{toc}{chapter}{Introduction to Part~I}
\label{ch:part1-intro}

In this Part, I present a comprehensive investigation of warm inflation (WI)~\cite{Berera:1995ie}, spanning theoretical foundations, computational tools, and novel observational signatures. WI offers a well-established alternative to the conventional inflationary picture --- referred to here as cold inflation (CI) --- in which the inflaton field is thermally coupled to a bath of radiation, continuously sourcing its production throughout the accelerated expansion. This is in contrast with the CI scenario, where interactions between the inflaton and other degrees of freedom are only relevant after inflation ends, during the so-called reheating phase, when the inflaton releases its energy via decay processes leading to the familiar radiation-dominated hot Big Bang cosmology. The continuous dissipation of the inflaton's energy into radiation introduces an additional friction term in the inflaton's equation of motion, parametrized by a dissipation rate $\Upsilon$ and conveniently expressed through the dimensionless dissipation strength $Q \equiv \Upsilon/(3H)$. These dissipative effects allow WI to naturally address several difficulties of the standard CI picture. The additional thermal friction relaxes the slow-roll conditions, allowing WI models to sustain a prolonged accelerated expansion for smaller, sub-Planckian, field excursions --- a regime in which the effective field theory description remains under better control and the tension with the Swampland conjectures in string theory is alleviated~\cite{Das:2018hqy,Motaharfar:2018zyb,Brandenberger:2020oav}. The dissipative coupling to the radiation bath affects not only the background evolution of the inflaton but also the nature of its perturbations,  which become predominantly thermal rather than quantum-mechanical in origin and fundamentally alter the predictions for CMB observables. In particular, the thermal contribution to the scalar power spectrum effectively suppresses the tensor-to-scalar ratio $r$, which allows WI to reconcile some of the simplest inflaton potentials, otherwise excluded by CMB data in the context of CI~\cite{Planck:2018jri,BICEP:2021xfz,Benetti:2016jhf}. Furthermore, the persistent radiation bath provides a smooth transition from the inflationary phase to radiation domination, alleviating the need for a separate reheating epoch and the theoretical uncertainties associated with it.

Despite these appealing features, the main challenge faced by WI since its inception has been to embed it in a consistent microphysical theory~\cite{Yokoyama:1998ju}. The core difficulty is that the same interactions responsible for generating the inflaton's dissipation also tend to produce thermal and radiative corrections to its potential, that are large enough to spoil the required flatness. Early constructions addressed this problem by introducing a large number of heavy intermediate fields in supersymmetric settings, whose large masses suppressed the dangerous thermal corrections to the inflaton sector~\cite{Berera:2008ar,Bastero-Gil:2012akf}. More recently, a different and arguably more economical approach has been developed, exploiting the shift symmetry of axion-like inflaton fields, which naturally protects the potential from large corrections while still allowing for efficient energy dissipation through the coupling to bosonic and fermionic fields~\cite{Bastero-Gil:2016qru, Berghaus:2019whh}. 

These developments triggered a renewed and sustained interest in WI and motivated a systematic exploration of its theoretical predictions and observational consequences. In this context, the first four chapters of this Part address both the theoretical foundations and the observational predictions of WI. On the theoretical side, I investigate the necessary conditions for a successful WI scenario. Specifically, in Chapter~\ref{ch:wi-scalar-field}, I derive general bounds on the required flatness of the inflaton potential in the strongly dissipative regime of WI, while in Chapter~\ref{ch:wi-eternal} I determine the conditions under which WI can avoid the onset of eternal inflation. On the observational side, I develop in Chapter~\ref{ch:warmspy} the public code \texttt{WarmSPy} for computing the full numerical primordial power spectrum in WI, ensuring accurate predictions for CMB observables across a broad range of dissipative regimes and microphysical scenarios. In Chapter~\ref{ch:wi-natural}, I examine how the WI framework performs when applied to natural inflation~\cite{Freese:1990rb} --- a model of particular interest for being one of the most economical and physically motivated inflationary constructions, and for its direct connection to the axion-like microphysical realizations of WI discussed above.

The final two chapters of this Part explore new observational signatures that are unique to the WI paradigm and arise generically in any scenario where a thermal bath persists throughout the primordial accelerated expansion. In Chapter~\ref{ch:wifi-dm}, I introduce a novel mechanism --- Warm Inflation Freeze-In (WIFI) --- through which the full dark matter abundance can be established during inflation itself, via feeble interactions with the persistent thermal bath. In Chapter~\ref{ch:thermal-gravitons}, I show that particle collisions in the WI bath inevitably generate an irreducible stochastic gravitational-wave background with a distinctive spectral shape, providing a complementary observational window into the thermal history of inflation.





